\documentclass[../main.tex]{subfiles}

\begin{document}

\section{Chapter 9: Storytelling and Presentation Narrative}
\label{sec:data-storytelling-design-opening}

This chapter reads the Vietnam Financial Inclusion Dashboard 2025 as a policy
brief built from weighted survey evidence rather than as a loose collection of
charts. Its central argument is that financial inclusion in Vietnam cannot be
judged by account ownership alone; it must be read as a progression from social
structure to digital readiness, from access to actual usage, and finally to
financial resilience.

\dashboardxref{fig:dashboard-overview}{Page 1 macro-socioeconomic baseline}

\subsection{The Core Policy Story}
\label{subsec:dashboard-purpose-storytelling-lens}

The chapter follows one narrative spine only:

\begin{center}
\textbf{Context \(\rightarrow\) Readiness \(\rightarrow\) Access \(\rightarrow\) Usage \(\rightarrow\) Resilience}
\end{center}

This sequence is persuasive because each page answers the most natural policy
question created by the page before it. A country can appear digitally
connected yet still fail to activate digital identity. It can appear broadly
banked yet still settle routine transactions in cash. It can appear financially
included yet still leave vulnerable households dependent on informal support
when shocks arrive.

\subsubsection{Page 1: Establishing the Structural Baseline}
\label{subsubsec:page1-context}

Page 1 serves as the structural baseline of the story, showing that later gaps
in access and usage emerge inside an already uneven social landscape rather than
in a neutral population.

The strongest evidence comes from the macro indicators and
\figref{fig:socioeconomic-matrix}{Socioeconomic matrix}. The page reports an Urban
Concentration Rate of only \(30.33\%\), an In Workforce Rate of \(87.92\%\),
and a Secondary Education Threshold Rate of \(89.98\%\). Yet the matrix shows
that the Poorest 20\% group has only 6 respondents with tertiary education or
more, compared with 72 in the Richest 20\% group. Those numbers make inequality
in income and education visible before the dashboard ever begins talking about
banking.

That opening matters for the policy argument because it prevents the reader from
treating later financial outcomes as purely technological problems. The story
begins by showing that exclusion is already structured by geography,
socioeconomic position, and educational opportunity, so any later digital or
financial intervention must be judged against that baseline rather than against
national averages alone.

\dashboardxref{fig:socioeconomic-matrix}{Socioeconomic matrix}

\subsubsection{Page 2: Readiness Exists, but Activation Is Weaker}
\label{subsubsec:page2-readiness}

Page 2 advances the story by testing whether the population that looks
demographically ready for participation is also ready to use digital financial
tools in practice.

\begin{figure}[H]
    \centering
    \begin{subfigure}[t]{0.31\textwidth}
        \centering
        \includegraphics[width=\textwidth]{fig-page2-connectivity-depth.png}
        \caption{Connectivity depth: 3M vs. 7D.}
        \label{fig:page2-connectivity-depth}
    \end{subfigure}
    \hfill
    \begin{subfigure}[t]{0.31\textwidth}
        \centering
        \includegraphics[width=\textwidth]{fig-page2-digital-vs-foundational-id.png}
        \caption{Digital vs. foundational ID coverage.}
        \label{fig:page2-id-coverage}
    \end{subfigure}
    \hfill
    \begin{subfigure}[t]{0.31\textwidth}
        \centering
        \includegraphics[width=\textwidth]{fig-page2-active-digital-id.png}
        \caption{Active digital ID verification rate.}
        \label{fig:page2-active-digital-id}
    \end{subfigure}
    \caption{Page 2 readiness evidence used in the storytelling analysis.}
    \label{fig:page2-readiness-overview}
\end{figure}

\dashboardfigure{0.86}{fig-page2-digital-activity-profile.png}
{Digital activity profile: penetration depth among internet users.}
{fig:page2-digital-activity-profile}

\figref{fig:page2-readiness-overview}{Page 2 readiness overview} shows that
internet access is both
high and frequent: Internet Use in the last three months reaches \(87.85\%\),
while Internet Use in the last seven days remains close at \(85.39\%\). The
same figure also shows a sharp activation gap: Foundational ID coverage reaches
\(99.65\%\), Online Digital ID reaches \(71.19\%\), but Active Digital ID
Verification falls to only \(32.68\%\). \figref{fig:page2-digital-activity-profile}{Digital activity profile}
pushes the argument further: Social Media reaches \(94.46\%\) and Read News
reaches \(78.44\%\), but Gov activity is only \(15.63\%\) and Job activity is
only \(4.89\%\).

Together, these figures argue that connectivity alone is not the same as
activation. Vietnam's adults appear highly present online, but their digital
behavior remains much stronger in communication and information than in
institutional or economic use. For a policy reader, this reframes the problem:
the bottleneck is not simply device or internet availability, but the
conversion of digital presence into verified, trusted, high-value use.

\subsubsection{Page 3: Access Is Broad, but Not Even}
\label{subsubsec:page3-access}

Page 3 becomes the diagnostic center of the story by asking who is formally
included, who is left behind, and what kind of inclusion the system is actually
delivering.

\dashboardfigure{0.82}{fig-page3-income-gender.png}
{Banking inclusion by income group and gender.}
{fig:page3-income-gender}

\figref{fig:account-kpi-cards}{Page 3 KPI cards} reports Banking Inclusion at
\(74.93\%\), Financial Institution Account ownership at \(74.05\%\), Mobile
Money Account ownership at \(46.69\%\), and Digital Account ownership at
\(67.03\%\). That already shows a 27.36 percentage-point gap between financial
institution accounts and mobile money accounts. \figref{fig:page3-income-gender}{Banking inclusion by income group and gender}
and \figref{fig:account-composition-flow}{Account composition flow} then show
that access is not socially neutral: lower-income groups, especially on the
rural side of the account-composition flow, are more exposed to unbanked or
thinner account states. \figref{fig:exclusion-friction-log}{Exclusion friction log}
adds a barrier layer to that story by showing that exclusion is not just a
missing account status but a set of recurring frictions tied to money, cost,
and access conditions.

\noindent{\small See \figref{fig:account-kpi-cards}{Page 3 KPI cards},
\figref{fig:page3-filter-bar}{Page 3 filter bar},
\figref{fig:account-composition-flow}{Account composition flow}, and
\figref{fig:exclusion-friction-log}{Exclusion friction log}.}

The policy consequence is that high national access cannot be read as uniform
inclusion. A dashboard that stopped at the \(74.93\%\) headline would suggest a
largely solved problem; this page instead shows that the remaining gap is
concentrated, structured, and diagnostically meaningful. The story therefore
shifts from ``how many people are banked?'' to ``which groups remain thinly
served, and what keeps them there?''

\subsubsection{Page 4: Access Does Not Guarantee Digital Usage}
\label{subsubsec:page4-usage}

Page 4 challenges the access story by asking whether accounts actually change
how money moves through everyday life.

\dashboardfigure{0.56}{fig-page4-digital-merchant-payments.png}
{Digital merchant and in-store payment by geography.}
{fig:page4-digital-merchant-payments}

\dashboardfigure{0.54}{fig-page4-utility-payment-methods.png}
{Utility payment methods by geography.}
{fig:page4-utility-payment-methods}

\figref{fig:national-inflows}{National inflow digitization} shows that
Corporate/Formal Wages are mostly
digitized, with \(80.99\%\) delivered directly to an account and only
\(18.60\%\) cash-only. Agricultural Sales show the opposite pattern:
\(73.42\%\) are cash-only and only \(25.79\%\) are delivered digitally to an
account. \figref{fig:page4-digital-merchant-payments}{Digital merchant payments by geography} adds a geographic gap:
Urban Merchant Payments reach \(65.79\%\), compared with \(50.30\%\) in rural
areas, while Urban In-Store Digital Pay reaches \(62.67\%\), compared with
\(46.11\%\) in rural areas. \figref{fig:page4-utility-payment-methods}{Utility payment methods by geography}
visually reinforces the same divide by showing that rural usage leans more
heavily toward cash-only or non-bill-paying states than the urban profile.

\dashboardxref{fig:national-inflows}{National inflow digitization}

This is the chapter's most important argumentative turn. Page 3 could have been
read as a success story about account ownership; Page 4 shows why that reading
is incomplete. The policy issue is not only whether people possess accounts,
but whether rural payment ecosystems, merchants, and routine household
transactions are digitized strongly enough for those accounts to matter in
practice.

\subsubsection{Page 5: The Real Test Is Resilience Under Stress}
\label{subsubsec:page5-resilience}

Page 5 gives the story its outcome test by asking whether participation in the
financial system is associated with stronger household resilience rather than
with access alone.

\dashboardfigure{0.78}{fig-page5-kpi-cards.png}
{Page 5 resilience KPI cards.}
{fig:page5-kpi-cards}

\dashboardfigure{0.68}{fig-page5-coping-income-group.png}
{Coping channels by income group.}
{fig:page5-coping-income-group}

\figref{fig:page5-kpi-cards}{Page 5 resilience KPI cards} shows a national
macro split in behavior:
\(73.05\%\) of respondents saved, while \(45.98\%\) borrowed. That national
picture is not the endpoint, however. \figref{fig:save-borrow-trajectories}{Save vs. borrow trajectories}
shows that saving and borrowing move differently across income groups and income
runway lengths, while \figref{fig:page5-coping-income-group}{Coping channels by income group}
reveals a much sharper difference in coping channels. In the Poorest 20\%
group, Family \& Friends stands at 86 while Formal Savings is only 30; in the
Richest 20\% group, Formal Savings rises to 60 while Family \& Friends falls
to 34. The full page overview in \figref{fig:page5-resilience-overview}{Page 5 resilience overview}
shows how these KPI,
trajectory, and coping views work together inside one resilience frame.

The final implication is that inclusion should be judged by how households
absorb shocks, not only by whether they hold an account. The dashboard ends on
resilience because it reveals a substantive difference between formal financial
participation and fallback dependence on informal networks. That ending gives
the chapter its human and policy-facing conclusion: the strongest financial
systems are not just those that open accounts, but those that reduce
vulnerability when households are under pressure.

\subsection{Why the Sequence Persuades}
\label{subsec:story-structure}

The sequence persuades because it prevents the reader from jumping too quickly
to a positive conclusion about inclusion. The chapter begins with structural
conditions, moves through readiness and access, tests whether those conditions
produce real digital usage, and only then asks whether that usage translates
into resilience. The same story can still be summarized once as
\textit{Context \(\rightarrow\) Readiness \(\rightarrow\) Access \(\rightarrow\)
Usage \(\rightarrow\) Resilience}, but the force of the sequence lies in what
each transition corrects.

The key corrective turn happens after Page 3. Banking Inclusion reaches
\(74.93\%\), and Financial Institution Account ownership reaches \(74.05\%\),
which could easily support a success narrative. Page 4 interrupts that reading
by showing Agricultural Sales at \(73.42\%\) cash-only and Urban Merchant
Payments at \(65.79\%\) versus \(50.30\%\) in rural areas. Page 5 then
interrupts any remaining optimism by showing that the Poorest 20\% still rely
far more on Family \& Friends at 86 than on Formal Savings at 30. In other
words, the chapter persuades because it steadily narrows the distance between
headline access and lived financial security.

\noindent{\small See \figref{fig:account-kpi-cards}{Page 3 KPI cards},
\figref{fig:national-inflows}{National inflow digitization}, and
\figref{fig:page5-coping-income-group}{Coping channels by income group}.}

\subsection{How the Evidence Builds the Argument}
\label{subsec:storytelling-principles}

The argument is built through accumulation rather than assertion. Each page
begins with a high-level anchor, then moves into a more diagnostic chart that
either qualifies or complicates the headline. For example, Page 2 does not stop
at high internet-use rates of \(87.85\%\) and \(85.39\%\); it contrasts them
with Active Digital ID Verification at only \(32.68\%\), then deepens the point
through \figref{fig:page2-digital-activity-profile}{Digital activity profile}, where Social Media
reaches \(94.46\%\) but Gov activity is only \(15.63\%\). That sequence turns a
connectivity story into an activation story.

The same evidence-building logic continues in later pages. Page 3 moves from
the \(74.93\%\) banking headline to social segmentation and exclusion barriers.
Page 4 moves from digital wage inflows at \(80.99\%\) to the much weaker rural
and agricultural usage environment. Page 5 moves from a national Saved rate of
\(73.05\%\) to the more revealing contrast between Family \& Friends at 86 and
Formal Savings at 30 for the Poorest 20\%. Because each visual layer narrows
the meaning of the previous one, the reader is led toward a more cautious and
policy-useful interpretation of inclusion.

\noindent{\small See \figref{fig:page2-readiness-overview}{Page 2 readiness overview},
\figref{fig:page2-digital-activity-profile}{Digital activity profile},
\figref{fig:account-kpi-cards}{Page 3 KPI cards},
\figref{fig:page4-utility-payment-methods}{Utility payment methods by geography},
and \figref{fig:page5-coping-income-group}{Coping channels by income group}.}

\subsection{Narrative Style and Claim Discipline}
\label{subsec:narrative-style}

The narrative style is persuasive, but it stays disciplined by grounding every
major interpretation in visible comparison rather than in speculative language.
The chapter does not claim that digital access \emph{causes} resilience or that
cash usage \emph{proves} exclusion. Instead, it repeatedly uses measured gaps to
show where the financial system appears stronger, where it appears weaker, and
which groups or settings deserve closer policy attention. That is why the tone
works better as a policy brief than as a promotional story.

This discipline matters most where the chapter could easily overstate its case.
The gap between Online Digital ID at \(71.19\%\) and Active Digital ID
Verification at \(32.68\%\) supports a claim about weaker activation, not a
claim that connectivity has failed altogether. The contrast between Formal
Wages at \(80.99\%\) direct-to-account and Agricultural Sales at \(73.42\%\)
cash-only supports a claim about uneven digitization, not a claim that all
rural finance is static. Likewise, the Poorest 20\% reliance on Family \&
Friends at 86 versus Formal Savings at 30 supports a claim about vulnerability,
not a deterministic claim about all low-income households. The chapter is most
credible when it keeps that distinction clear.

\noindent{\small See \figref{fig:page2-digital-activity-profile}{Digital activity profile},
\figref{fig:national-inflows}{National inflow digitization},
\figref{fig:page4-digital-merchant-payments}{Digital merchant payments by geography},
and \figref{fig:page5-kpi-cards}{Page 5 resilience KPI cards}.}

\subsection{Interaction as Policy Exploration}
\label{subsec:interaction-exploration-story}

The chapter's policy argument is strengthened by interaction because the
dashboard does not force the reader to stop at the national averages used in
the narrative. Page 3 begins with Banking Inclusion at \(74.93\%\), Financial
Institution Account ownership at \(74.05\%\), Mobile Money Account ownership at
\(46.69\%\), and Digital Account ownership at \(67.03\%\), then gives the user
filters for Urbanicity, Income Group, Gender, Education Level, and In
Workforce. Page 5 does something similar from a resilience baseline of
\(73.05\%\) saved and \(45.98\%\) borrowed, using filters for Location, Gender,
and In Workforce to move the reader from national averages into subgroup
patterns.

This interaction design matters because the chapter's most important claims are
distributional. The national story says that digital activation is weaker than
connectivity, that account access is broader than digital usage, and that
resilience differs across social groups. The filters on Page 3 and Page 5 make
those claims investigable instead of merely rhetorical. At the same time, the
available source is still a static PDF export, so the report should only claim
that the dashboard is \emph{designed for} subgroup exploration, not that every
filtered result has been directly tested from the exported file.

\noindent{\small See \figref{fig:page3-filter-bar}{Page 3 filter bar},
\figref{fig:page5-filter-bar}{Page 5 filter controls},
\figref{fig:page3-interaction-states}{Page 3 default account-access state}, and
\figref{fig:page5-resilience-overview}{Page 5 resilience overview}.}

\end{document}
